\documentclass[10pt,a4paper]{article}
\usepackage[margin=20mm,headheight=15pt,headsep=7mm,footskip=10mm]{geometry}
\usepackage{fontspec,amsmath,amssymb,mathtools,microtype}
\setmainfont{lmroman10-regular.otf}[BoldFont=lmroman10-bold.otf,ItalicFont=lmroman10-italic.otf,BoldItalicFont=lmroman10-bolditalic.otf]
\setsansfont{lmsans10-regular.otf}[BoldFont=lmsans10-bold.otf,ItalicFont=lmsans10-oblique.otf,BoldItalicFont=lmsans10-boldoblique.otf]
\setmonofont[Scale=.86]{lmmono10-regular.otf}[ItalicFont=lmmono10-italic.otf]
\usepackage{ragged2e}\usepackage{xcolor,booktabs,longtable,array,graphicx,enumitem,fancyhdr,needspace,titlesec,hyperref,xurl}
\definecolor{AtlasNavy}{HTML}{18334C}\definecolor{AtlasTeal}{HTML}{116D73}\definecolor{AtlasGray}{HTML}{52616D}
\hypersetup{colorlinks=true,linkcolor=AtlasTeal,urlcolor=AtlasTeal,pdfcreator={Complexity Atlas},unicode=true}
\urlstyle{same}\setlength{\parindent}{0pt}\setlength{\parskip}{4.8pt}\setlength{\emergencystretch}{3em}
\renewcommand{\baselinestretch}{1.08}\setcounter{secnumdepth}{0}\setlist{nosep,leftmargin=1.5em}
\titleformat{\section}{\Large\bfseries\sffamily\color{AtlasNavy}\raggedright}{}{0em}{}
\titleformat{\subsection}{\large\bfseries\sffamily\color{AtlasTeal}\raggedright}{}{0em}{}
\titleformat{\subsubsection}{\normalsize\bfseries\color{AtlasNavy}\raggedright}{}{0em}{}
\titlespacing*{\section}{0pt}{14pt}{6pt}\titlespacing*{\subsection}{0pt}{10pt}{5pt}
\pagestyle{fancy}\fancyhf{}\fancyhead[L]{\small\sffamily\color{AtlasGray}COMPLEXITY ATLAS}
\fancyhead[R]{\small\sffamily\color{AtlasGray}ENGLISH EDITION}\fancyfoot[C]{\small\thepage}\renewcommand{\headrulewidth}{.3pt}
\newcommand{\AtlasMath}[1]{\par\begingroup\centering\setbox0=\hbox{$\displaystyle #1$}\ifdim\wd0>1.4\linewidth\typeout{ATLAS WIDE FORMULA: \the\wd0}\fi\ifdim\wd0>\linewidth\resizebox{\linewidth}{!}{\box0}\else\box0\fi\par\endgroup}
\newcommand{\AtlasInlineMath}[1]{\begingroup\setbox0=\hbox{$#1$}\ifdim\wd0>.4\linewidth\endgroup\AtlasMath{#1}\else\endgroup$#1$\fi}
\pdfstringdefDisableCommands{\def\AtlasInlineMath#1{#1}}
\DeclareMathOperator{\polylog}{polylog}\DeclareMathOperator{\dist}{dist}\DeclareMathOperator{\diam}{diam}
\DeclareMathOperator{\rank}{rank}\DeclareMathOperator{\tr}{tr}\DeclareMathOperator{\diag}{diag}
\DeclareMathOperator{\prox}{prox}\DeclareMathOperator{\KL}{KL}\DeclareMathOperator{\Var}{Var}
\providecommand{\eps}{\varepsilon}\providecommand{\R}{\mathbb R}\providecommand{\E}{\mathbb E}
\providecommand{\mathbbm}[1]{\mathbf{#1}}\providecommand{\mathds}[1]{\mathbf{#1}}
\providecommand{\llbracket}{[\![}\providecommand{\rrbracket}{]\!]}
\begin{document}
\raggedright

{\LARGE\bfseries\sffamily\color{AtlasNavy} Finite-sum nonconvex optimization: selectable component-value queries\par}\vspace{5pt}
\textbf{Setting ID.} \texttt{SO\_\allowbreak{}ZO\_\allowbreak{}C05\_\allowbreak{}rand\_\allowbreak{}expected}

Single-objective optimization

Find an approximate first-order stationary point of the population objective \AtlasInlineMath{f}: an output point with sufficiently small gradient norm. The mathematical object is \AtlasInlineMath{\nabla f}, and information is acquired through the zeroth-order interface specified below.

\textbf{Status.} Polynomial gap present

\textbf{TASK tags.} formulation: First order stationarity; solution concept: First order stationarity; task family: Single objective

\textbf{MODEL tags.} convexity: Convexity not assumed; initialization: Global initial gap; smoothness: Componentwise gradient lipschitz; structure: Finite sum

\textbf{ORACLE tags.} access order: Zeroth-order interface; interface family: Selected component value; public information: Public inputs explicitly declared; query domain: Full space

\textbf{ALGORITHMS tags.} adaptivity: Adaptive; randomization: Private randomness allowed; restriction: No additional span restriction explicitly imposed

\textbf{OUTPUT tags.} certificate: No additional output certificate required

\textbf{METRIC tags.} criterion: First order stationarity

\textbf{GUARANTEE tags.} aggregation: Per instance; mode: Expectation

\textbf{RESOURCE tags.} budget: Fixed worst-case budget

\needspace{7\baselineskip}\section{1. Setting specification}
\needspace{4\baselineskip}\subsection{Task}
Find an approximate first-order stationary point of the population objective \AtlasInlineMath{f}: an output point with sufficiently small gradient norm. The mathematical object is \AtlasInlineMath{\nabla f}, and information is acquired through the zeroth-order interface specified below.

\needspace{4\baselineskip}\subsection{Admissible model}
\AtlasInlineMath{f:\mathbb{R}^d\to\mathbb{R}} is differentiable and satisfies \AtlasInlineMath{\|\nabla f(x)-\nabla f(y)\|_2\le\beta\|x-y\|_2} for every \AtlasInlineMath{x,y\in\mathbb{R}^d}.

\AtlasInlineMath{f(x)=n^{-1}\sum_{i=1}^n f_i(x)}, where each \AtlasInlineMath{f_i:\mathbb{R}^d\to\mathbb{R}} is \AtlasInlineMath{\beta}-smooth. Individual components may be nonconvex.

\AtlasInlineMath{f_*:=\inf_{x\in\mathbb{R}^d}f(x)>-\infty}, and the given initial point \AtlasInlineMath{x_0} satisfies \AtlasInlineMath{f(x_0)-f_*\le\Delta}.

\needspace{4\baselineskip}\subsection{Oracle interface}
The algorithm may actively choose \AtlasInlineMath{i\in\{1,\ldots,n\}} and \AtlasInlineMath{x\in\mathbb{R}^d}. The interface \AtlasInlineMath{(i,x)\mapsto f_i(x)} returns one exact component value. Components may be revisited; full-average values must be computed through component calls.

The public inputs are \AtlasInlineMath{d}, \AtlasInlineMath{\beta}, \AtlasInlineMath{\Delta}, \AtlasInlineMath{\varepsilon}, \AtlasInlineMath{n}, \AtlasInlineMath{x_0}, together with the sets, kernels and auxiliary operations explicitly declared above. Information about the unknown objective and any undisclosed sample distribution is obtained only through the stated interface.

\needspace{4\baselineskip}\subsection{Output}
Return \AtlasInlineMath{\widehat{x}\in \mathbb{R}^d}. The output may be an unqueried point, a random choice among previous points, or an average or convex combination that remains in the stated space. No function value, gradient, inner solution, or accuracy certificate must accompany the output.

\needspace{4\baselineskip}\subsection{Loss or accuracy criterion}
The base loss is \AtlasInlineMath{\ell(I,x)=\|\nabla f(x)\|_2} with tolerance \AtlasInlineMath{\varepsilon>0}. Any theorem for squared gradient norms must first be converted to this norm scale.

\needspace{4\baselineskip}\subsection{Resource accounting}
One \AtlasInlineMath{(i,x)\mapsto f_i(x)} costs \AtlasInlineMath{1} component-value calls; recomputing the full average costs \AtlasInlineMath{n} calls. Every initialization, snapshot and differencing point is charged individually.

Only the declared oracle resource is charged. Computation and storage using existing information, together with explicitly known geometric and proximal operations, are free.

\needspace{4\baselineskip}\subsection{Algorithm views and success guarantees}
Private randomization is allowed, with the seed independent of the instance and external randomness. Queries, stopping and postprocessing may be arbitrary measurable adaptive functions of the history. No linear-span or zero-respecting restriction is imposed.

Interactions are sequential, and every actual response in a batch is charged separately. Except for explicitly listed auxiliary interfaces, arithmetic, linear algebra, storage, and internal computation on available information are free, with no additional communication or memory restriction. Finite-difference steps must be finite and positive; real-number precision is not charged separately.

After fixing the public inputs, choose one algorithm that works for every admissible instance and reply law. Require \AtlasInlineMath{\mathbb{E}[\ell(I,\widehat{x})]\le \varepsilon}, fixing the instance before averaging over algorithmic and feedback randomness.

Oracle cost uses a deterministic worst-case budget: \AtlasInlineMath{N\le T} on every possible execution path. For multiple resources, coordinatewise caps apply jointly along the same path.

\needspace{7\baselineskip}\section{2. Notation and parameters}
\AtlasInlineMath{d}: Dimension of the variable space \textbf{Domain:} \AtlasInlineMath{d\in\mathbb{N},\ d\ge1}.

\AtlasInlineMath{\beta}: Public upper bound on the gradient Lipschitz constant \textbf{Domain:} \AtlasInlineMath{\beta>0}.

The object, norm and domain are fixed by the specific inequalities in this profile's model.

\AtlasInlineMath{\Delta}: Known upper bound on the initial objective gap \textbf{Domain:} \AtlasInlineMath{\Delta\ge0}.

\AtlasInlineMath{f(x_0)-\inf_{x\in\mathbb{R}^d}f(x)\le\Delta}

\AtlasInlineMath{\varepsilon}: Error target; norm and squared-norm criteria use different thresholds \textbf{Domain:} \AtlasInlineMath{\varepsilon>0}.

\AtlasInlineMath{n}: Number of finite-sum components \textbf{Domain:} \AtlasInlineMath{n\in\mathbb{N},\ n\ge1}.

\AtlasInlineMath{F=n^{-1}\sum_{i=1}^n f_i}

\AtlasInlineMath{K}: Dimensionless iteration scale for smooth stationarity \textbf{Domain:} \AtlasInlineMath{K\ge0}.

\AtlasInlineMath{K=\beta\Delta/\varepsilon^2}

\AtlasInlineMath{x,y,u}: Function arguments and dummy variables in the formulas \textbf{Domain:} \AtlasInlineMath{x,y,u\in\mathbb{R}^d}; spherical domains are embedded in \AtlasInlineMath{\mathbb{R}^{d+1}}.

Admissible ranges are specified separately by the task, model and oracle in this profile.

\AtlasInlineMath{\widehat{x}}: Final algorithm output \textbf{Domain:} \AtlasInlineMath{\mathbb{R}^d}.

\AtlasInlineMath{I}: Complete problem instance \textbf{Domain:} \AtlasInlineMath{I\in\mathcal{I}}.

Includes the unknown objective and its admissible feedback distribution or environment law. Public parameters are fixed before the quantifiers.

\AtlasInlineMath{\mathcal{I}}: Admissible instance class determined by the eight specification fields in this profile \textbf{Domain:} Instances satisfying the public parameter bounds.

\AtlasInlineMath{\mathcal{A}}: Algorithm or strategy \textbf{Domain:} Measurable adaptive strategies allowed in this profile.

The algorithm is selected before the hidden instance.

\AtlasInlineMath{\ell}: Base loss or residual \textbf{Domain:} \AtlasInlineMath{\ell:\mathcal{I}\times\mathcal{X}_{\mathrm{out}}\to[0,+\infty]}.

The exact formula appears in this profile's performance-criterion section.

\AtlasInlineMath{\mathbb{E},\Pr}: Expectation and probability \textbf{Domain:} Evaluated over the algorithmic and feedback randomness specified by this profile's guarantee.

Fix the instance first; aggregate cost and error separately.

\AtlasInlineMath{\inf,\sup}: Infimum and supremum \textbf{Domain:} Evaluated only over the variable sets explicitly specified in the formula.

In the risk definition, the infimum over algorithms precedes the supremum over instances or centers.

\AtlasInlineMath{\|\cdot\|_2,\langle\cdot,\cdot\rangle}: Euclidean norm and inner product \textbf{Domain:} \AtlasInlineMath{\|x\|_2=\sqrt{\langle x,x\rangle}}.

\AtlasInlineMath{\langle x,y\rangle=x^\top y}

\AtlasInlineMath{B_p(c,R)}: Closed norm ball \textbf{Domain:} \AtlasInlineMath{p\in[1,\infty],\ R\ge0}.

\AtlasInlineMath{B_p(c,R)=\{x:\|x-c\|_p\le R\}}; \AtlasInlineMath{B_2} denotes the Euclidean ball.

\AtlasInlineMath{N,T}: Actual call count and declared budget \textbf{Domain:} \AtlasInlineMath{N,T\in\mathbb{N}_0}.

A fixed budget bounds every execution path; the Markov specification uses an expected budget.

\AtlasInlineMath{O,\Omega,\Theta,\widetilde{O}}: Asymptotic upper bounds, lower bounds, matching orders, and upper bounds with hidden logarithms \textbf{Domain:} Within the joint parameter window specified for each complexity regime.

Dependence of hidden constants on fixed parameters is stated separately for each bound. Tilde notation may hide logarithmic factors only in the declared variables.

\AtlasInlineMath{c,C,c_\alpha,C_\alpha}: Positive constant in a proof or asymptotic bound \textbf{Domain:} \AtlasInlineMath{0<c,C,c_\alpha,C_\alpha<\infty}.

Unsubscripted constants have the dependencies stated in the bound; constants subscripted by \AtlasInlineMath{\alpha} may depend only on the fixed smoothness order.

\AtlasInlineMath{\log,\log_+,e}: Natural logarithm, truncated logarithm, and Euler's number \textbf{Domain:} \AtlasInlineMath{z>0}.

\AtlasInlineMath{\log_+(z)=\max\{0,\log z\}}; \AtlasInlineMath{e=\exp(1)}, and \AtlasInlineMath{\log_2} denotes the base-two logarithm.

\AtlasInlineMath{f_i,i}: Finite-sum components and their indices \textbf{Domain:} \AtlasInlineMath{f_i:\mathbb{R}^d\to\mathbb{R},\ i\in\{1,\ldots,n\}}.

Components are unknown; indices may be selected actively and revisited at a cost.

\AtlasInlineMath{F}: Finite-sum mean objective \textbf{Domain:} \AtlasInlineMath{F:\mathbb{R}^d\to\mathbb{R}}.

\AtlasInlineMath{F=n^{-1}\sum_{i=1}^n f_i}

\AtlasInlineMath{f,f_*}: Finite-sum objective and its global infimum \textbf{Domain:} \AtlasInlineMath{f:\mathbb{R}^d\to\mathbb{R}}.

\AtlasInlineMath{f=F}, \AtlasInlineMath{f_*=\inf_{x\in\mathbb{R}^d}f(x)>-\infty}.

\AtlasInlineMath{\nabla f,\nabla^2f,D^jf}: Gradient, Hessian, and higher Fréchet derivatives \textbf{Domain:} \AtlasInlineMath{\nabla f\in\mathbb{R}^d,\ \nabla^2f\in\mathbb{R}^{d\times d}}; \AtlasInlineMath{D^jf} is a \AtlasInlineMath{j}-linear form.

Used only at derivative orders and on regions where the model guarantees its existence.

\AtlasInlineMath{\|\cdot\|_{\mathrm{op}}}: Euclidean induced matrix or tensor operator norm \textbf{Domain:} Take the supremum over Euclidean unit input vectors of the corresponding order.

\AtlasInlineMath{\|T\|_{\mathrm{op}}=\sup_{\|u_1\|_2,\ldots,\|u_j\|_2\le1}|T[u_1,\ldots,u_j]|}

\AtlasInlineMath{\mathcal{X}_{\mathrm{fun}},\mathcal{X}_{\mathrm{out}},\mathcal{X}}: Function domain, admissible output domain and optimization domain \textbf{Domain:} Each is the Euclidean space or subset explicitly specified in this profile.

\AtlasInlineMath{\mathcal{X}_{\mathrm{out}}=\mathbb{R}^d} and \AtlasInlineMath{\mathcal{X}=\mathbb{R}^d}; the function domain is the region on which the model imposes regularity.

\AtlasInlineMath{N^*}: Optimal worst-case oracle complexity satisfying the complete guarantee \textbf{Domain:} \AtlasInlineMath{N^*\in[0,+\infty]}.

Take the infimum over admissible algorithms, followed by the worst case over the instances of this profile. With multiple resources, use this scalar only for an explicitly specified resource projection.

\AtlasInlineMath{x_0}: Public initial point \textbf{Domain:} \AtlasInlineMath{x_0\in\mathbb{R}^d}.

Belongs to the feasible domain required by the model.

\needspace{12\baselineskip}\section{3. Complexity bracket}
\needspace{4\baselineskip}\subsection{Compared quantities}
\needspace{4\baselineskip}\subsubsection{Q1 - target accuracy cost}
\AtlasMath{C^\star_{v,\pi}(\varepsilon;\theta)=\inf\{B\ge0:\exists\mathcal A\in\mathfrak A_v,\ C_\pi(\mathcal A,I)\le B\ \text{and}\ G_v(\mathcal A,I;\varepsilon)\ \forall I\}}
\textbf{Definition.} One scalar component value \AtlasInlineMath{f_i(x)}

\textbf{Scope.} declared\_joint\_parameter\_window

\textbf{Optimization direction.} minimize

Here v is the selected algorithm view, G\_v its complete success guarantee, I an admissible fixed instance, theta the remaining parameters, and C\_pi the declared resource projection. C\_pi uses the view's fixed-path or expected-work budget convention, including initialization. No change of guarantee or scalarization is implicit.

\needspace{4\baselineskip}\subsection{R1 - Parameter regime C05.R1}
\textbf{Region.} Set \AtlasInlineMath{K=\beta\Delta/\varepsilon^2}; require \AtlasInlineMath{K\ge C}, \AtlasInlineMath{n\ge1024} and \AtlasInlineMath{\varepsilon^2\le c\beta\Delta}. The dimension satisfies both \AtlasInlineMath{d\ge CK\log(2+K)} and \AtlasInlineMath{d\ge C(1+K/\sqrt n)\log(2+n+\sqrt nK)}.

\textbf{Quantity.} Q1

\textbf{Resource.} One scalar component value \AtlasInlineMath{f_i(x)}

\textbf{Status.} Polynomial gap

\textbf{Lower bound.} B2, B3

\AtlasMath{\Omega(\max\{dK,n+\sqrt n K\})}
\textbf{Upper bound.} B1

\AtlasMath{O(d(n+\sqrt n K))}
\textbf{Gap or comparison.} \AtlasInlineMath{O(d(n+\sqrt nK)/\max\{dK,n+\sqrt nK\})}

\needspace{7\baselineskip}\section{4. Bounds and supporting evidence}
The following conditions apply to every bound in this profile, in addition to each bound's stated region and algorithm scope.

\textbf{Oracle mapping.} Each source component value corresponds to one selectable component query; computing the exact mean requires accessing every component.

\textbf{Guarantee mapping.} Error uses this profile's accuracy scale and aggregation order. The source theorem's expectation, probability, and asymptotic quantifiers are retained separately.

\textbf{Output mapping.} The output need only lie in this setting's prescribed space. Permitted averages and unqueried outputs must still satisfy its feasibility conditions.

\textbf{Resource mapping.} Atomic calls under a fixed cap on every execution path. All feedback obtained during initialization, differencing and inner iterations is charged.

\needspace{4\baselineskip}\subsection{B1 - Zeroth-order SPIDER variance reduction for finite sums}
\textbf{Admission.} Included in the bracket

\textbf{Direction.} upper

\textbf{Witness or lower-bound scope.} Randomized witness

\AtlasMath{O\!\left(d[n+\sqrt n\,\beta \Delta/\varepsilon^2]\right)}
\textbf{Sources.} \href{https://arxiv.org/pdf/1807.01695}{[2]}

\textbf{Conditions.} Every component has the stated gradient smoothness, and the initial objective gap is at most \AtlasInlineMath{\Delta}. A complete snapshot costs on the order of \AtlasInlineMath{dn} scalar values and is included in the total budget.

\textbf{Source location.} \href{https://arxiv.org/pdf/1807.01695}{[2]} theorem: Finite-sum version of Theorem 8, including initialization costs; pdf pages: 20; printed pages: 20

\needspace{4\baselineskip}\subsection{B2 - Weakly hidden finite-sum hard family}
\textbf{Admission.} Included in the bracket

\textbf{Direction.} lower

\textbf{Witness or lower-bound scope.} All randomized adaptive algorithms

\AtlasMath{\Omega(d\beta \Delta/\varepsilon^2)}
\textbf{Sources.} \href{https://arxiv.org/html/2610.19652v1}{[3]} \href{https://arxiv.org/html/2609.00045v3}{[1]} \href{https://arxiv.org/pdf/2610.00245v2}{[4]} \href{https://arxiv.org/html/2610.17074v1}{[5]}

\textbf{Conditions.} Use the exact-value hard subclass with identical components, retaining the accuracy and high-dimensional conditions of the corresponding nonconvex lower bound.

\textbf{Local argument.} A1

\textbf{Source location.} \href{https://arxiv.org/html/2610.19652v1}{[3]} theorem: Theorem 10.5, with all components identical

\textbf{Source location.} \href{https://arxiv.org/html/2610.19652v1}{[3]} theorem: Theorem 10.5, with all components identical; pdf pages: 24; printed pages: 24

\needspace{4\baselineskip}\subsection{B3 - Weakly hidden finite-sum hard family}
\textbf{Admission.} Included in the bracket

\textbf{Direction.} lower

\textbf{Witness or lower-bound scope.} All randomized adaptive algorithms

\AtlasMath{\Omega(n+\sqrt n\,\beta \Delta/\varepsilon^2)}
\textbf{Sources.} \href{https://arxiv.org/html/2609.00045v3}{[1]} \href{https://arxiv.org/html/2610.19652v1}{[3]} \href{https://arxiv.org/pdf/2610.00245v2}{[4]} \href{https://arxiv.org/html/2610.17074v1}{[5]}

\textbf{Conditions.} \AtlasInlineMath{n\ge1024,\varepsilon^2\le c\beta \Delta}; set \AtlasInlineMath{K=\beta \Delta/\varepsilon^2}, for which \AtlasInlineMath{d\ge C(1+K/\sqrt n)\log(2+n+\sqrt nK)} is sufficient. The incremental first-order oracle returns \AtlasInlineMath{(f_i,\nabla f_i)}.

\textbf{Local argument.} A1

\textbf{Source location.} \href{https://arxiv.org/html/2609.00045v3}{[1]} theorem: Theorem 3.1; Section 3 equation (4)

\textbf{Source location.} \href{https://arxiv.org/html/2609.00045v3}{[1]} theorem: Theorem 3.1; Section 3 equation (4); pdf pages: 5; printed pages: 5

\needspace{7\baselineskip}\section{5. Local arguments}
\needspace{4\baselineskip}\subsection{A1 - Oracle and subclass transfers for nonconvex lower bounds}
\textbf{Scope.} Each bound keeps the source's fixed cap, output convention, accuracy window, and lower-bound algorithm class. Dimension-free supremum statements remain distinct from fixed-dimension statements.

The registered lower bounds follow by exact simulation or class containment, with no added dimension factor.

For identical stochastic samples or identical finite-sum components, use the same exact-value hard function for every label. The oracle response is then exactly one source value, so the source's computational lower bound transfers one-for-one. When the target observes fresh additive Gaussian values, an exact-value algorithm can generate the independent Gaussian perturbations itself; an algorithm for that noisier interface would therefore contradict the stronger exact-value lower. The exact-value source cap \AtlasInlineMath{B=\min\{\beta\Delta,L^2\}} and its high-dimensional window remain in force. A smooth-only class may choose the proof-only budget \AtlasInlineMath{L=\sqrt{\beta\Delta}} inside the hard subclass; this does not add a public Lipschitz parameter.

The second-moment sample class is contained in each \AtlasInlineMath{r}-moment class for \AtlasInlineMath{1<r\le2}, by \AtlasInlineMath{(\mathbb EG^r)^{1/r}\le(\mathbb EG^2)^{1/2}}. Thus the source Goldstein statistical lower transfers to the larger heavy-tail class, but it does not establish a matching heavy-tail exponent. A composite class that includes \AtlasInlineMath{\psi=0} contains the corresponding noncomposite hard family. With \AtlasInlineMath{X=\mathbb R^d}, the normal-cone residual is the unconstrained Goldstein residual. These embeddings support the registered composite and constrained-normal-cone lower bounds only where those model members are allowed; they do not automatically support a compact-set Frank--Wolfe-gap lower bound.

A value-and-gradient incremental oracle can simulate a selected-component value by discarding the gradient. Dense26 Theorem 3.1 consequently transfers its \AtlasInlineMath{\Omega(n+\sqrt n\,\beta\Delta/\varepsilon^2)} lower, with \AtlasInlineMath{n\ge1024} and the registered dimension window. The same stronger-oracle ordering transfers the Hessian-Lipschitz and higher-order-smooth first-order lower theorems to exact values. A second-order stationary output necessarily satisfies its first-order gradient requirement, so the first-order lower also limits that stronger task. If the theorem states that some sufficiently large finite dimension is hard, retain \AtlasInlineMath{\sup_{d\ge1}N^*(d,\varepsilon)} or its explicit hard-dimension bound; do not present it as a result for every fixed \AtlasInlineMath{d}.

\textbf{Limits of the argument.} Local argument and specification transfer; cited external theorems remain primary dependencies.

\textbf{Sources.} \href{https://arxiv.org/html/2610.19652v1}{[3]} \href{https://arxiv.org/html/2609.00045v3}{[1]} \href{https://arxiv.org/pdf/2610.00245v2}{[4]} \href{https://arxiv.org/html/2610.17074v1}{[5]}

\needspace{7\baselineskip}\section{References}
[1] Peng, Tang and Jia. Dense Weak Hiding: Closing Complexity Gaps in Nonconvex and PL Finite-Sum Optimization under Individual Smoothness. arXiv:2609.00045v3, 2026-10-10.

\url{https://arxiv.org/html/2609.00045v3}

[2] Fang, Li, Lin and Zhang. SPIDER: Near-Optimal Non-Convex Optimization via Stochastic Path-Integrated Differential Estimator. NeurIPS 2018; arXiv:1807.01695v2.

\url{https://arxiv.org/pdf/1807.01695}

[3] Xia and Yuan. On Global Joint Lower Bounds in Zeroth-Order Nonconvex Optimization. arXiv:2610.19652v1, 2026-10-14.

\url{https://arxiv.org/html/2610.19652v1}

[4] Zhang et al.. Matching Lower Bounds for Randomized First-Order Methods in Hessian-Lipschitz Nonconvex Optimization. arXiv:2610.00245v2, 2026-10-19.

\url{https://arxiv.org/pdf/2610.00245v2}

[5] Zhang et al.. Matching Randomized First-Order Bounds for Higher-Order Smooth Nonconvex Optimization. arXiv:2610.17074v1, 2026.

\url{https://arxiv.org/html/2610.17074v1}

\end{document}
